% !TeX root = ../papers/nc-p3.tex
\subsection{三次元接点の三生成基底が滑らかな平面三次曲線を持つ場合}
\label{sec:elliptic-tangent-windows}
三生成基底 $B=T(H)/(S)$ の三つの双一次関係が
滑らかな平面三次曲線 $E\subset\mathbb P(H^*)$ 上の自己同型 $\sigma$ のグラフを
定める枝を扱う。以下 $L=\mathcal O_E(1)$ とする。

\begin{theorem}\label{thm:elliptic-tangent-classification}
この枝で有限四辺を持つ標準AS窓の基底は、平面三次曲線 $E$ と、
その上の位数 $m=2$ または $4$ の平行移動 $\sigma$ から得るものに限り、全て実現する。
各基底に対する全補正は定理\ref{thm:ternary-ore-linear-space}の九次元空間 $P_B$ である。
非零の全ての補正が有限四辺を与え、零のみ辺が正次元である。
非零な射影補正は次数12の有効因子 $D$ として復元できる。
二つの窓の添字保持同型は、$(E,L,\sigma,D)$ を対応させる曲線同型の存在と同値である。
最小周期は $\sigma^jD$ を $D$ と同じ $\Aut(E,L,\sigma)$ 軌道に送る最小の $j>0$ である。
二つの位数の族はそれぞれ既約で同型類の次元は9、一般最小周期は $m$ である。
\end{theorem}

\begin{proof}[滑らかな曲線上での四移動の必要性]
定理\ref{thm:maximal-tangent-ore}によって $B$ は三次元標準二次ASである。
唯一の三次交叉を張るテンソル $w$ の先頭・末端の縮約を
$\ell,r:H^*\xrightarrow{\sim}S$ とする。連続する三点を $p,q=\sigma p,u=\sigma^2p$
とし、混合補正 $d_i$ の $(p,q)$ での評価を $v_i\in H^*$ とする。
$p,q$ が独立なら、二つの第四座標が存在するのは $v_i\in\langle p,q\rangle$ の場合だけである。
$\sigma$ が恒等なら、三次曲線には消える二次式がないため $S$ は交代式の全空間となり、
そのグラフの台は平面全体になって矛盾する。従って一般に $p,q$ は独立である。
有限辺なら一般の $p$ で $v_i\notin\langle p,q\rangle$ である。

$A_u(v)=(s(v,u))_{s\in S}$、$B_p(v)=(s(p,v))_{s\in S}$ と書く。
両写像の核は $kq$ で、像は $S^*$ の平面である。
式\eqref{eq:ternary-shift-derivation}を $(p,q,u)$ で評価すれば
$A_u(v_i)+B_p(v_{i+1})=0$ となる。
$A_u(p)=B_p(u)\ne0$（一般に $u\ne q$）で、$A_u(v_i)$ はこれと独立である。
よって二平面は等しい。その零化直線はそれぞれ $kr(\sigma u)$ と
$k\ell(\sigma^{-1}p)$ である。実際 $w(-,u,\sigma u)=w(\sigma^{-1}p,p,-)=0$ である。
従って全曲線上で $\sigma^4=\mathbb P(r^{-1}\ell)|_E$ である。

もし $\sigma p=p$ なら $r(p)$ と $\ell(p)$ は比例する。
線型行列 $B_p$ の階数は2で、その核は $kp$、余核の双対は $k\ell(p)$ である。
従って行列式の微分は $v\mapsto\ell(p)(v,p)$ の定数倍だが、
これは $r(p)(v,p)=w(v,p,p)=0$ に比例する。行列式は $E$ を定める三次式なので
滑らかさに反する。従って固定点はなく、$\sigma$ は非恒等な平行移動である。
実際、楕円曲線の自己同型の原点保持部分 $\alpha$ が非恒等なら $1-\alpha$ は全射である。
\end{proof}

\begin{proof}[基底の表示と全ファイバーの正則性]
変曲点 $O$ を原点にしたHesse表示
$x^3+y^3+z^3-3hxyz=0$、$O=[1:-1:0]$ を使う。
この表示は標数零の代数閉体上の滑らかな平面三次曲線を覆う。
例えば変曲点を原点にWeierstrass表示にすれば埋込みは $|3O|$ であり、
Hesse族の $j=27h^3(h^3+8)^3/(h^3-1)^3$ は射影直線から全射で
特異パラメータは極だけである\cite[\S7.3]{KLV}。
$\sigma$ の平行移動点 $P=[a:b:c]$ は三等分点でない。
三等分点なら $\sigma$ は平面に延び、二次式の制限の単射性から
$S$ は射影変換の平面グラフの三関係となる。従って $abc\ne0$ である。
ここで座標軸との九交点は全て変曲点であり、接線を代入しても確かめられる。

文字を $x_0,x_1,x_2$ として、三つの関係を巡回順に
\[
          c x_j^2+a x_{j+1}x_{j+2}+b x_{j+2}x_{j+1}
\]
とする。左点 $p=(x,y,z)$ に対する行列の三行は
$(cx,bz,ay),(az,cy,bx),(by,ax,cz)$ である。
行列式は $(a^3+b^3+c^3)xyz-abc(x^3+y^3+z^3)$ で、$E$ 上の全点で階数2である。
逆方向も $a,b$ を交換すれば同じで、グラフを定める。
次点 $q=(v_0,v_1,v_2)$ に対し、最初の二関係の差は
$p,P,(v_1,v_0,v_2)$ の共線条件である。
最初の二座標の交換は $O$ に関する群逆（$q$ と交換像を結ぶ直線は $O$ を通る）なので、
このグラフは $P$ の平行移動を与える。
グラフで消える双一次式の空間は三次元である。
実際、$L\not\simeq\sigma^*L$ なので、$L$ の基点自由な二切断の完全列を
$\sigma^*L$ で捩ると、次数零の非自明な束の $H^1$ の消滅により
二次の積は次数6の全切断に全射となる。従って元の三関係もこの表示になる。
三関係には巡回対称な三次交叉 $w$（各関係に対応する文字を左から掛けて足したもの）があり、
$\ell=r$ である。先の必要条件は $4P=0$ となり、残る位数は2と4である。

逆にこれらの $P$ の全てで三生成基底が正則となることを示す。
$L_j=(\sigma^j)^*L$、$L_{[j,n)}=\bigotimes_{j\le h<n}L_h$ と置く。
$R_n=H^0(L_{[0,n)})$ に積 $f\cdot g=f(\sigma^{\deg f})^*g$ を入れる。
$n\ge3$ で、$B$ の生成元・関係・三次交叉による右複体を曲線上で評価した列
\[
 0\longrightarrow L_{[3,n)}\longrightarrow S\otimes L_{[2,n)}
  \longrightarrow H\otimes L_{[1,n)}\longrightarrow L_{[0,n)}\longrightarrow0
\]
は束として完全である。最後の写像は $p$ での評価、中央写像は第二文字に
$\sigma p$ を代入したもの、最初の写像は $w$ の末尾に $\sigma^2p$ を代入したものである。
各点で合成は零で、階数が順に $1,2,1$ なので完全となる。
$n>3$ では全項の $H^1$ が零で、大域切断の列も完全である。
$n=3$ では左端の $H^1(\mathcal O_E)=k$ だけが残り、
大域切断では $H\otimes H^0(L_{[1,3)})$ の位置に限り一次元のホモロジーがある。
$n=2$ での積の全射性と核は上の三関係の計算である。
従って $R$ は一次生成で $B\to R$ は全射である。

$k$ の右 $B$ 自由解消の最初の二微分は生成元と $S$ で与えられる。
その第三項に入れる全てのsyzygyに先の $w$ も含められる。
従ってこの解消を $R$ でテンソルしたときの第1ホモロジーは今計算した
$k(-3)$ であり、第2ホモロジーは零である。後者は $w$ の像だけですでに核を埋めるためである。
次数付き左 $B$ 加群 $R$ の極小解消から、$B\to R$ の核は
$B(-3)$ と同型で、右からある三次元 $g$ を掛けて埋め込まれる。
よって $H_B(t)=H_R(t)/(1-t^3)=(1-t)^{-3}$ である。
$R$ は整域であるので、$B$ の左socle元は $b'g$ と書け、右からの $g$ の単射性により
$b'$ も左socleに入る。次数降下で零となる。反対側も同じ証明である。
$w$ の両端支持が全三次元であるので、定理\ref{thm:maximal-tangent-ore}の末尾と同様に、
三次までのKoszul複体の最左写像は単射で、Hilbert級数から全列が完全となる。
完全な相補積から三次元左右AS条件も従う。
\end{proof}

\begin{proof}[九次元空間の一様な計算]
$a=1$ とし、対角変数変更と各関係のスカラー倍で三つの平方係数を $1,1,u$、
$u=c^3$ にできる。実際、各平方係数は変数倍率 $\lambda_j$ により
$c\lambda_j^3/(\lambda_0\lambda_1\lambda_2)$ となる。
位数2なら $a=b=1$、$u=t$ と置く。
位数4では $b=t$、$u=t(1+t^2)/(1+t)$ となる。
実際、$2P$ の座標は行列の二行の外積から
$(b(c^3-a^3),a(b^3-c^3),c(a^3-b^3))$ である。
三座標の三乗が全て等しい場合は特異なので除いてある。
$2P$ が $O$ でない二等分点となる条件は前二座標の一致であり、表示式を与える。
各場合 $bu\ne0$、$(1+b^3+u)^3\ne27b^3u$ を課し、位数4では分母 $1+t\ne0$ とする。
$t=1$ や $t^2+t+1=0$ は位数4の滑らかな領域に入らない。

$\omega$ を原始三乗根とする。$x_j\mapsto\omega^j x_j$ の三つの重みに分けて
$\lambda,\rho$ を計算できる。三つの重みの計算は、平方係数を一定に戻した基底での
文字の巡回置換により対応する。実際、二操作の交換子は文字全てへの同じ三乗根倍であり、
次数1の補正空間上で重みを一つ移す。
重み0では各 $x_j$ の二次補正に、$i+k=j\bmod3$, $i\ne k$ の二単項式 $x_ix_k$
を使う。平方は関係で消えるので六変数である。$\lambda,\rho$ は各々 $10\times6$ となる。
有理関数体上の厳密計算は次表となる。
\[
\begin{array}{c|c|c|c}
m&[\lambda\ \rho]\text{ の六次小行列式}
     &\mu&\rk(\lambda+\mu\rho)\\ \hline
2&-t(t-1)&1,-1&4,3\\
4&-t^2(t-1)(t^2+1)(t^2+t+1)/(t+1)^2&-1,\sqrt{-1},-\sqrt{-1}&4,4,4
\end{array}
\]
\path{computations/verify_elliptic_tangent_windows.py} は各三次関係の重みの
九つの三次語から $(S\otimes H)+(H\otimes S)$ の行を零化する汎関数を計算し、
その行に各補正を入れてこの表を再現する。小行列式に使う汎関数の分母の最小公倍式は
それぞれ $1,t(t+1)$ である。従って示した小行列式は全ての許される $t$ で有効かつ非零である。
一方、三次商が定数階数なので最後の欄の一般階数は全特殊値での上界である。

$\rho$ の核は右乗法のこの重みの一方向なので階数5である。
連結行列の階数下界6から、次項を許す初期値の次元は高々5である。
左右乗法の二方向は独立である。比例すれば $x_0 h$ と $h x_0$ が全ての文字で比例するが、
非零の平方係数を持つ三関係を比較すると $x_0x_1$ と $x_1x_0$ は独立である。
ゆえに商は一重み当たり高々3次元である。
$\ker(\lambda+\mu\rho)$ の元は $d_i=\mu^i d$ で両側に延び、移動の $\mu$ 固有類を与える。
補生成元の変更の二方向のうちこの核に入るものは $L_v-\mu R_v$ の一直線だけである。
実際 $\lambda R_v=\rho L_v$ は零でなく（$\ker\rho=R_H$）、$\lambda L_v=\rho R_v=0$ である。
よって表から $m=2$ の固有次元は少なくとも $1,2$、
$m=4$ ではそれぞれ少なくとも1となる。異なる固有値は独立なので全て等号であり、
三重みを合わせて常に九次元となる。特に検算に使った有理点への特殊化だけに依存しない。
\end{proof}

\begin{proof}[有限条件、同型、特殊化と周期]
$N=L\otimes\sigma^*L$ と置く。$B_2=H^0(N)$ なので、補正 $d$ の評価ベクトル
$v(p)$ は $H^*\otimes H^0(N)$ に属する。
$p,\sigma p$ が全点で独立であることから
\[
 0\longrightarrow \sigma^*L\oplus L\longrightarrow H^*\otimes N
 \xrightarrow{\ \det(p,\sigma p,-)\ } N^2\longrightarrow0
\]
は完全である。最初の写像は二つの係数を $p,\sigma p$ に掛けて足す。
$H^1$ の消滅と一次の完全線型系から、大域切断での核はちょうど左右乗法である。
従って $P_B$ は $H^0(N^2)$ 内の九次元部分空間と同一視される。
非零切断の零因子 $D$ は長さ12である。各零点では独立な二行から第四座標二つを解けるため、
持ち上がった局所環は曲線上の因子の局所環と一致する。
一方だけで基底への射影が零なら他方でも零となる。
残る指定点は有限完全交叉の全長20から長さ8で、唯一の三次元接点である。
特殊値で重根が生じれば重複度がそのまま残余局所長になる。
切断零では全曲線が持ち上がる。

一意な点と三生成部分を保つ同型は $(E,L,\sigma)$ の射影同型を誘導し、
逆にこの同型はグラフで消える三関係を対応させる。
同じ因子を与える非零切断の比は定数であり、この定数は文字全てを同じ定数倍する
基底自己同型で実現できる。従って因子を含む表示が全ての特殊値でも必要十分である。
$\Aut(E,L)$ は有限である。例えば九つの変曲点の像を試し、四点の射影枠で変換を定めて
曲線・グラフ・因子を比較できる。隣接辺のスキームが同じ点列を与えるので、
次の残余因子は重複度を込めて $\sigma D$ となる。

各位数の基底パラメータは上の $t$ の非空開部分で尽くされる。
その像は一次元である。Hesseの $j$ は各場合に非定数であり、
固定した楕円曲線では変曲原点と所定の平行移動点はいずれも有限個である。
補正の核と商は上の定数階数計算によりベクトル束をなし、射影ファイバーは八次元、
残る射影群は有限なので像は既約九次元である。
固定基底について、九次元切断空間の線型系の像は $\mathbb P^8$ を張る曲線で次数8以上である。
固定零因子を除いても移動部分の次数は高々12なので、この写像は像へ双有理である。
もし一般因子の $j$ 移動を有限群 $\Aut(E,L,\sigma)$ で戻せれば、
群のある一元が切断の全射影空間でこの移動と一致する
（各候補の一致は射影線型な閉条件）。その元と $\sigma^j$ の差は全因子を保ち、
上記の双有理写像により恒等である。平行移動 $\sigma^j$ が $L$ を保つには
$3jP=0$ が必要十分で、$m=2,4$ では $m\mid j$ である。これが一般最小周期を与える。
位数2と4は滑らかな領域内では互いに特殊化しない。変曲原点を局所的に選んだ後の
四等分点の群スキームは有限エタールで、二等分点はその開かつ閉な部分である。
\end{proof}

定理\ref{thm:algebraic-nine}の九変数族の一般点では、三生成部分の曲線は滑らかで、
移動は固定点のない対合であり、指定点は接空間三次元である。
従ってその九次元の像は上の位数2の九次元族に稠密である。
今回の因子による同型判定は、この共通部分で定理\ref{thm:algebraic-nine}の
次元と一般周期の計算を補う。全モジュライでの既約成分性はこの比較からは出ない。
